%/*% pdflatex -file-line-error -interaction nonstopmode %

\documentclass{article}
\usepackage{graphicx}
\graphicspath{ {./imagescopy/standard_test/} }
\usepackage{amsmath}
\usepackage{amsfonts}
\usepackage{hyperref}
\usepackage{exercises}
\usepackage{amsthm}
\usepackage{yhmath}
\usepackage{placeins}
\newcommand*{\oldneg}{\mathord{\sim}}

\theoremstyle{definition}
\newtheorem{definition}{Definition}

\setlength{\parindent}{.5in}
\setlength{\parskip}{12pt}

\title{2.1 Answers}
\begin{document}
\maketitle

\section*{1}
(a): Product of Powers Property.
\newline
(b): Negative Exponent Property.
\newline
(c): Power of a Product Property.

\section*{2}
Scientific notation is defined to be a number in the form $c\times 10^n$ where $1\le c\le 10$ and $n$ is an integer. The number $25.2\times 10^{-3}$ is 
such that $c=25.2>10$, thus meaning the given number is not in scientific notation.

\section*{3}
$3^3\times 3^2=3^{3+2}=3^5$ (Prodcut of Powers Property)

\section*{4}
$(4^{-2})^3=4^{-2\times 3}=4^{-6}$ (Power of a Power Property)
$\newline =\frac{1}{4^6}$ (Negative Exponent Property)

\section*{5}
$(-5)(-5)^4=(-5)^1\times (-5)^4=(-5)^5$ (Product of Powers Property)

\section*{6}
$(2^4)^2=2^{4\times 2}=2^8$ (Power of a Power Property)

\section*{7}
$\frac{5^2}{5^5}=5^{2-5}=5^{-3}$ (Quotient of Powers Property)
$\newline =\frac{1}{5^3}$ (Negative Exponent Property)

\section*{8}
$(\frac{3}{5})^4=\frac{3^4}{5^4}=\frac{81}{625}$ (Power of a Quotient Property)

\section*{9}
$(\frac{2}{7})^{-3}=\frac{2^{-3}}{7^{-3}}$ (Power of a Quotient Property)
$\newline =\frac{\frac{1}{8}}{\frac{1}{343}}$ (Negative Exponent Property)
$\newline =\frac{343}{8}$

\section*{10}
$9^3\times 9^{-1}=9^{3-1}=9^2$ (Product of Powers Property)

\section*{11}
$\frac{3^4}{3^{-2}}=3^{4+2}=3^6$ (Quotient of Powers Property)

\section*{12}
$(\frac{2}{3})^{-5}(\frac{2}{3})^4=(\frac{2}{3})^{-5+4}=(\frac{2}{3})^{-1}$ (Product of Powers Property)
$\newline =\frac{1}{\frac{2}{3}}=\frac{3}{2}$ (Negative Exponents Property)

\section*{13}
$6^3\times 6^0\times 6^{-5}=6^{3+0}\times 6^{-5}=6^3\times 6^{-5}=6^{3-5}=6^-2$ (Product of Powers Property)
$\newline =\frac{1}{36}$ (Negative Exponents Property)

\section*{14}
$((\frac{1}{2})^{-5})^2=(\frac{1}{2})^{-5\times 2}=(\frac{1}{2})^{-10}$ (Power of a Power Property)
$\newline =\frac{1}{(\frac{1}{2})^{10}}$ (Negative Exponents Property)
$\newline =\frac{1}{\frac{1^{10}}{2^{10}}}=\frac{1}{\frac{1}{1024}}=1024$ (Power of a Quotien Property)

\section*{15}
$(4.2\times 10^3)(1.5\times 10^6)=(4.2\times 1.5)(10^3\times 10^6)=6.3\times 10^9$

\section*{16}
$(1.2\times 10^{-3})(6.7\times 10^{-7})=(1.2\times 6.7)(10^{-3}\times 10^{-7})=8.04\times 10^{-10}$

\section*{17}
$(6.3\times 10^5)(8.9\times 10^{-12})=(6.3\times 8.9)(10^5\times 10^{-12})=56.07\times 10^{-7}=5.607\times 10^1\times 10^{-7}=5.607\times 10^{-6}$

\section*{18}
$(7.2\times 10^9)(9.4\times 10^8)=(7.2\times 9.4)(10^9\times 10^8)=67.68\times 10^{17}=6.768\times 10^{18}$

\section*{19}
$(2.1\times 10^{-4})^3=2.1^3\times (10^{-4})^3=9.261\times 10^{-12}$

\section*{20}
$(4.0\times 10^3)^4=4.0^4\times (10^3)^4=256\times 10^{12}=2.56\times 10^{14}$

\section*{21}
$\frac{8.1\times 10^{12}}{5.4\times 10^9}=\frac{8.1}{5.4}\times \frac{10^{12}}{10^9}=1.5\times 10^3$

\section*{22}
$\frac{1.1\times 10^{-3}}{5.5\times 10^{-8}}=\frac{1.1}{5.5}\times \frac{10^{-3}}{10^{-8}}=0.2\times 10^5=2\times 10^4$

\section*{23}
$\frac{(7.5\times 10^8)(4.5\times 10^{-4})}{1.5\times 10^7}=\frac{(7.5\times 4.5)(10^8\times 10^{-4)}}{1.5\times 10^7}\newline\rightarrow 
=\frac{3.375\times 10^5}{1.5\times 10^7}=2.25\times 10^{-2}$

\section*{24}
$\frac{w^{-2}}{w^6}=w^{-2-6}=w^{-8}$ (Quotient of Powers Property)
$\newline =\frac{1}{w^8}$ (Negative Exponents Property)

\section*{25}
$(2^2y^3)^5=4^5\times (y^3)^5$ (Power of a Product Property)
$\newline =1024y^{15}$ (Power of a Power property)

\section*{26}
$(p^3q^2)^{-1}=(p^3)^{-1}\times (q^2)^{-1}$ (Power of a Product Property)
$\newline =p^{-3}q^{-2}$ (Power of a Power Property)
$\newline =\frac{1}{p^3q^2}$ (Negative Exponents Property)

\section*{27}
$(w^3x^{-2})(w^6x^{-1})=(w^3\times w^6)(x^{-2}\times x^{-1})=w^9x^{-3}$ (Product of Powers Property)
$\newline =\frac{w^9}{x^3}$ (Negative Exponents Property)

\section*{28}
$(5s^{-2}t^4)^{-3}=(5s^{-2})^{-3}\times (t^4)^{-3}$ (Power of a Product Property)
$\newline =5^{-3}s^6t^{-12}$ (Power of a Product Property, Power of a Power Property)
$\newline =\frac{s^6}{125t^{12}}$ (Negative Exponents Property)

\section*{29}
$(3a^3b^5)^{-3}=(3a^3)^{-3}\times (b^5)^{-3}$ Power of a Product Property)
$\newline =3^{-3}a^{-9}b^{-15}$ (Power of a Product Property, Power of a Power Property)
$\newline =\frac{1}{27a^9b^{15}}$ (Negative Exponents Property)

\section*{30}
$\frac{x^{-1}y^2}{x^2y^{-1}}=x^{-1-2}y^{2+1}=x^{-3}y^3$ (Quotient of a Power Property)
$\newline =\frac{y^3}{x^3}$ (Negative Exponent Property)

\section*{31}
$\frac{3c^3d}{9cd^{-1}}=\frac{c^2d^2}{3}$ (Quotient of Powers Property)

\section*{32}
$\frac{4r^4s^5}{24r^4s^{-5}}=\frac{r^0s^{10}}{6}=\frac{s^{10}}{6}$ (Quotient of Powers Property)

\section*{33}
$\frac{2a^3b^{-4}}{3a^5b^{-2}}=\frac{2a^{-2}b^{-2}}{3}$ (Quotient of Powers Property)
$\newline =\frac{2}{3a^2b^2}$ (Negative Exponents Property)

\section*{34}
\[
\begin{array}{rl}
            & \frac{y^{11}}{4z^3}\times\frac{8z^7}{y^7} \\
\rightarrow & \frac{8z^7y^{11}}{4z^3y^7} \\
\rightarrow & 2z^4y^4\ Quotient\ of\ Powers\ Property
\end{array}
\]

\section*{35}
\[
\begin{array}{rl}
            & \frac{x^2y^{-3}}{3y^2}\times \frac{y^2}{x^{-4}} \\
\rightarrow & \frac{x^2y^{-3}y^2}{x^{-4}3y^2} \\
\rightarrow & \frac{x^2y^{-1}}{x^{-4}3y^2}\ Product\ of\ Powers\ Property \\
\rightarrow & \frac{x^6y^{-3}}{3}\ Quotient\ of\ Powers\ Property \\
\rightarrow & \frac{x^6}{3y^3}\ Negative\ Exponents\ Property
\end{array}
\]

\section*{36}
B

\section*{37}
The error here is that the powers of the denominator and numerator were divided. The Quotient of Powers Property states that the powers should in 
fact be subtracted, meaning the correct answer is $x^{10-2}=x^8$

\section*{38}
The error here is that the powers of the factors of the product were multiplied. The Product of Powers Property actually states that the powers should be 
added, not multiplied, meaning the correct answer is $x^{5+3}=x^8$.

\section*{39}
The error here is that $(-3)$ and $(-3)$ were multiplied together. However, there is no properrty that $a^m b^n=(ab)^{m+n}$ where $m\ne n$, meaning 
this conclusion is incorrect. Instead, we can use the Product of Powers property to conclude that $(-3)^2(-3)^4=(-3)^6$

\section*{40}
We are given that $A=\frac{\sqrt{3}}{4}s^2$. With the given value for $s$, we can find $A$ using this formula:
\[
\begin{array}{rl}
            & A=\frac{\sqrt{3}}{4}(\frac{x}{3})^2 \\
\rightarrow & \frac{\sqrt{3}}{4}\frac{x^2}{9} \\
\rightarrow & \frac{x^2\sqrt{3}}{36}
\end{array}
\]

\section*{41}
\[
\begin{array}{rl}
            & V=\pi r^2h \\
\rightarrow & \pi x^2\frac{x}{2} \\
\rightarrow & \frac{\pi x^3}{2}
\end{array}
\]

\section*{42}
\[
\begin{array}{rl}
            & V=lwh \\
\rightarrow & \frac{5x}{3}2x\times x \\
\rightarrow & \frac{10x^3}{3}
\end{array}
\]

\section*{43}
\[
\begin{array}{rl}
            & x^{15}y^{12}z^8=x^4y^7z^{11}\times x^{11}y^5z^{-3} \\
\end{array}
\]

\section*{44}
\[
\begin{array}{rl}
            & 3x^3y^2=\frac{12x^2y^5}{4x^{-1}y^3} \\
\end{array}
\]

\section*{45}
\[
\begin{array}{rl}
            & (a^5b^4)^2=a^{14}b^{-1}\times a^{-4}b^9 \\
\end{array}
\]

\section*{46}
\[
\begin{array}{rl}
            & x^{12}y^{16}=(x^8y^8)(x^4y^8)=(x^6y^9)(x^6y^7)=(x^2y^{10})(x^{10}y^6) \\
\end{array}
\]

\section*{47}
\[
\begin{array}{rl}
            & a^{-n}=\frac{a^0}{a^n}\  Quotient\ of\ Powers\ Property \\
\rightarrow & \frac{a^0}{a^n}=\frac{1}{a^n}\  Zero\ Exponenet\ Property \\
\rightarrow & a^{-n}=\frac{1}{a^n}\  Transitive\ Property\ of\ Equality \\
\end{array}
\]

\section*{48}
\[
\begin{array}{rl}
            & (\frac{a}{b})^n=(a\times\frac{1}{b})^n=a^n(\frac{1}{b})^n\  Product\ of\ Powers\ Property \\
\rightarrow & a^n(\frac{1}{b})^n=a^n(b^{-1})^n\  Negative\ Exponents\ Property \\
\rightarrow & a^n(b^{-1})^n=a^nb^{-n}\ Power\ of\ a\ Power\ Property \\
\rightarrow & a^nb^{-n}=\frac{a^n}{b^n}\ Negative\ Exponents\ Property \\
\rightarrow & (\frac{a}{b})^n=\frac{a^n}{b^n}\ Transitive\ Property

\end{array}
\]

\section*{49}
Volume of Pacific $=(1.56\times 10^{14})(4.03\times 10^3)=(1.56\times 4.03)(10^{14}\times 10^3)=6.2868\times 10^{17}$\newline
Volume of Atlantic $=(7.68\times 10^{13})(3.93\times 10^3)=(7.68\times 3.93)(10^{13}\times 10^3)=30.1824\times 10^{16}=3.01824\times 10^{17}$\newline
Volume of Indian $=(6.86\times 10^{13})(3.96\times 10^3)=(6.86\times 3.96)(10^{13}\times 10^3)=27.1656\times 10^{16}=2.71656\times 10^{17}$\newline
Volume of Arctic $=(1.41\times 10^{13})(1.21\times 10^3)=(1.41\times 1.21)(10^{13}\times 10^3)=1.7061\times 10^{16}$

\section*{50}
We are given that the South American continent has been moving about $0.000022$ miles per year for the past $125,000,000$ years, meaning 
the total amount that the continent has moved in that time is $0.000022\times 125,000,000$. This can be simplified further into scientific notation:\[
\begin{array}{rl}
            & 0.000022\times 125,000,000 \\
\rightarrow & 0.000022\times 1.25\times 10^8 \\
\rightarrow & 0.0000275\times 10^8=2.75\times 10^3
\end{array}
\]

\section*{51}
We are comparing the volumes of a bead and a pearl, both of which we will assume are in spherical shapes, meaning their volume can be given by: 
$\frac{4}{3}\pi r^3$. The bead has a diameter of $6$ millimeters, or a radius of $3$ millimeters, while the pearl has a radius of $4.5$ millimeters, 
meaning we can compare their volumes:
\[
\begin{array}{rl}
            & \frac{\frac{4}{3}\pi 3^3}{\frac{4}{3}\pi 4.5^3} \\
\rightarrow & \frac{27}{91.125}=\frac{8}{27} 
\end{array}
\]
meaning the ratio of the volume of the bead to the volume of the pearl is $\frac{8}{27}$

\section*{52}
(a): Because we are considering three spheres, the total volume of the three tennis balls is $3(\frac{4}{3}\pi r^3)=4\pi r^3$\newline
(b): The volume of a cylinder (and this cylinder) is $\pi r^2h$\newline
(c): We know that the diameter of each of the tennis balls is $2r$. With the information that $h$ is equal to the sum of the three diameters of the tennis 
balls, we can conlude that $h=2r+2r+2r=6r$\newline
(d): Now that we have an expression of $h$ in terms of $r$, we can rewrite the volume of the can, and find how much of it is taken up by the three 
tennis balls:
\[
\begin{array}{rl}
            & \frac{4\pi r^3}{\pi r^26r} \\
\rightarrow & \frac{2}{3} 
\end{array}
\]
meaning the balls take up $\frac{2}{3}$ of the can.

\section*{53}
(a): $V=\pi 9.53^2\times 1.55=\pi \times 140.772395\ Cubic\ millimeters=1.40772395\pi \times 10^{-11}$ cubic meter.\newline
(b): ?\newline
(c): ?

\section*{54}
(a): We are given that the radius of earth is $5$ times that of the inner core, which has a radius of $r$, meaning the radius of earth is $5r$, 
and its volume is $\frac{4}{3}\pi (5r)^3=\frac{4}{3}\pi 125r^3$. The volume of the inner core is $\frac{4}{3}\pi r^3$, meaning the volume of Earth's total 
volume to the volume of the inner core is $\frac{\frac{4}{3}\pi 125r^3}{\frac{4}{3}\pi r^3}=125$\newline
(b): We are given that the radius of the outer core is $\frac{11}{6}r+r=\frac{17}{6}r$, meaning the volume of the outer core is $\frac{4}{3}\pi (\frac{17}{6}
r)^3=\frac{4}{3}\pi \frac{4913}{216}r^3$, meaning the ratio of the volume of Earth's outer core to the volume the inner core is:
$\frac{\frac{4}{3}\pi \frac{4913}{216}r^3}{\frac{4}{3}\pi r^3}=\frac{4913}{216}\approx 22.7$



\end{document}
